\chapter{The Descartes Universe}
\label{chap:descartes}

\section{The Descartes condition}
\label{sec:descartes-condition}

% ERROR: block "descartes-condition" not found
\begin{remark}[The Descartes condition on examples]
  \label{rem:descartes-calabi-yau}
  \uses{ex:calabi-yau, ex:ale}
The constraint $\chi(E) = -6$ specialises to each archimedean example:

\begin{itemize}
  \item \emph{Calabi--Yau} (\hyperref[ex:calabi-yau]{Example}): constrains
\end{itemize}
  the alternating sum of Betti numbers; only manifolds admitting
  a principal $\mathrm{SU}(2)$-bundle with this Euler characteristic
  realise the Descartes universe.
\begin{itemize}
  \item \emph{ALE links} (\hyperref[ex:ale]{Example}): selects finite subgroups
\end{itemize}
  $\Gamma \subset \mathrm{SU}(2)$ via arithmetic constraints on
  the McKay quiver.
\begin{itemize}
  \item \emph{Affine space} $M = \mathbb{C}^2$: forces non-trivial
\end{itemize}
  bundle topology; the $q$-harmonic zero-modes of $\Delta_q$
  reproduce the CMB thermal spectrum in the $q \to 1$ regime.
\end{remark}

\begin{remark}[Descent to one complex dimension]
  \label{rem:descent-one-dim}
  \uses{prop:jet-prolongation, prop:euler-char-stability, rem:descent-reduction-zero}
By the descent theorem (\hyperref[prop:jet-prolongation]{Proposition}), a
two-complex-dimensional Descartes universe $(M, G, P, E)$ with $M = \mathbb{C}^2$
descends to a one-complex-dimensional variety $M_0$.  The stability of the
Euler characteristic under descent and lifting
(\hyperref[prop:euler-char-stability]{Proposition}) ensures that the descended
universe retains $\chi(E) = -6$.  We may therefore study the Descartes
universe in the setting of one-complex-dimensional varieties.
The descent may continue further to dimension zero
(\hyperref[rem:descent-reduction-zero]{Remark}), where the entire quantum
observable universe is captured by fibre data at a single point.
\end{remark}

\begin{proposition}[Genus of the one-complex-dimensional descent]
  \label{prop:genus-maximal-descent}
  \uses{def:descartes-condition}
Let $(M, G, P, E)$ be a Descartes universe that descends to a compact
one-complex-dimensional variety $M_0$.  Then $M_0$ is a Riemann surface of
genus $g = 4$.

---

The Descartes condition gives $\mathrm{tr}(\mathrm{id}_A) = -6$.  Under the fibre functor $\omega$ of the archimedean realisation, this maps to the sheaf Euler characteristic $\chi(V) = -6$, where $V = P \times_{\mathrm{SU}(2)} K^2$ is the associated rank-$2$ vector bundle over $M_0$.  Since $G = \mathrm{SU}(2)$, the determinant bundle $\det(V)$ is trivial, so $\deg(V) = 0$.  By Riemann--Roch on the compact Riemann surface $M_0$: $\chi(V) = \deg(V) + \mathrm{rank}(V)(1-g) = 2(1-g)$.  Setting $2(1-g) = -6$ gives $g = 4$.
\end{proposition}
\begin{proof}
The Descartes condition gives $\mathrm{tr}(\mathrm{id}_A) = -6$.  Under the fibre functor $\omega$ of the archimedean realisation, this maps to the sheaf Euler characteristic $\chi(V) = -6$, where $V = P \times_{\mathrm{SU}(2)} K^2$ is the associated rank-$2$ vector bundle over $M_0$.  Since $G = \mathrm{SU}(2)$, the determinant bundle $\det(V)$ is trivial, so $\deg(V) = 0$.  By Riemann--Roch on the compact Riemann surface $M_0$: $\chi(V) = \deg(V) + \mathrm{rank}(V)(1-g) = 2(1-g)$.  Setting $2(1-g) = -6$ gives $g = 4$.
\end{proof}

\section{Leptons and chirality}
\label{sec:descartes-leptons-chirality}

\begin{remark}[Electrons, muons and the Reeb decomposition]
  \label{rem:descartes-leptons}
  \uses{def:reeb-decomposition, def:fermion}
The Reeb decomposition (\hyperref[def:reeb-decomposition]{Definition}) applied to
a Descartes universe yields two classes of fermionic excitations
(\hyperref[def:fermion]{Definition}).  The degree of the vertical component
$\alpha_V$ distinguishes them: forms with $\deg\alpha_V = 1$ correspond to
electrons, while forms with $\deg\alpha_V = 2$ correspond to muons.  More
generally, each additional unit of vertical degree is identified with a
successive lepton generation.  The constraint $\chi(E) = -6$ bounds the
number of such generations, since the Hodge decomposition of $E$ admits only
finitely many independent vertical harmonic forms.  This
generation-dependent structure leads to testable predictions, including a
$q$-correction to the proton charge radius that differs between muonic and
electronic hydrogen (\hyperref[ssec:proton-radius]{Section}).
\end{remark}

\begin{remark}[Chirality as Reeb orientation]
  \label{rem:descartes-chirality}
  \uses{def:fermion, def:boson}
The Reeb vector field $\xi$ on the CR manifold $M$ carries a natural
orientation: the sign of the contact form $\theta(\xi) = 1$ singles out a
preferred direction along the vertical distribution.  Left-handed and
right-handed chiralities correspond to the two orientations of $\xi$ relative
to the complex structure of $M$.  For fermions
(\hyperref[def:fermion]{Definition}), the vertical component $\alpha_V$ transforms
with a definite sign under reversal of $\xi$, so chirality is encoded in the
$\pm$-eigenspaces of the Reeb involution $\xi \mapsto -\xi$.  Bosons
(\hyperref[def:boson]{Definition}), being purely horizontal, are chirality-blind ---
consistent with the observed parity invariance of the electromagnetic and
strong interactions.
\end{remark}

\begin{remark}[Light cone of an affine chart]
  \label{rem:light-cone-affine}
Consider an open subset $U \subset M_0$ representing a portion of the
descended universe from the big bang to heat death.  An \emph{affine chart}
on $U$ is a local coordinate patch in which the metric on the base manifold
$M_0$ defines a causal structure.  The \emph{light cone} of the affine chart
is the boundary of the causal future in $U$: it is determined by the metric
on $M_0$ and has the property that no singularity of the ambient space is
visible from within the chart.  In this way, the affine chart provides a
non-singular window onto a portion of the Descartes universe.
\end{remark}

\section{Trefoil subgroups and Bring's curve}
\label{sec:trefoil-brings}

\begin{remark}[Chirality and trefoil subgroups]
  \label{rem:chirality-trefoil}
  \uses{prop:genus-maximal-descent}
The chirality of the observable universe --- manifest in the three colour
charges of quantum chromodynamics --- is encoded in the symmetry group of the
descended surface $M_0$.  \hyperref[prop:genus-maximal-descent]{Proposition}
identifies $M_0$ as a genus-$4$ Riemann surface.  We require a genus-$4$
surface whose automorphism group contains three subgroups isomorphic to $S_3$
--- quotients of the trefoil knot group
$\pi_1(S^3 \setminus K) \cong \langle x, y \mid x^2 = y^3 \rangle$ ---
corresponding to the three colour charges (red, green, blue).  These $S_3$
subgroups generate a subgroup of the monodromy group, and the remaining
generators of the full automorphism group act by conjugation, permuting the
three $S_3$ subgroups among themselves.
\end{remark}

\begin{remark}[Bring's curve]
  \label{rem:brings-curve}
  \uses{rem:chirality-trefoil}
Bring's curve is the unique Riemann surface of genus $4$ with the maximal
automorphism group $S_5$ (order $120$) for its genus.  It is defined by
the intersection of three symmetric hypersurfaces in $\mathbb{P}^4$:

\[
\sum_{i=1}^{5} z_i = 0,
  \qquad
  \sum_{i=1}^{5} z_i^2 = 0,
  \qquad
  \sum_{i=1}^{5} z_i^3 = 0.
\]

The symmetric group $S_5$ acts by permuting the five homogeneous coordinates,
and this action realises the monodromy structure described in
\hyperref[rem:chirality-trefoil]{Remark}: the group $S_5$ contains three copies of
$S_3$ (quotients of the trefoil knot group) corresponding to the three
colour charges, and elements outside these subgroups act by conjugation,
permuting the three copies among themselves.
\end{remark}

\section{Colour-sector coordinates on the exceptional fibre}
\label{sec:colour-sector-coords}

% ERROR: block "exceptional-divisor" not found
% ERROR: block "chiral-sector-coords" not found
\section{Categorical braid structure}
\label{sec:descartes-braid}

\begin{corollary}[Braid-group structure of a Descartes universe]
  \label{cor:descartes-braid-structure}
  \uses{def:descartes-condition, prop:braid-action-tensor, def:su2-generators, rem:chirality-trefoil, def:chiral-sector-coords, rem:r-matrix-braiding, ex:trefoil-operator-algebra, prop:brings-is-qou, ex:p-adic-realisation}
Let $(\mathcal{C}, A, \omega)$ be a Descartes universe
(\hyperref[def:descartes-condition]{Definition}).  Then:

\begin{enumerate}
  \item The braiding of $\mathcal{C} = \mathrm{Rep}(\mathrm{SU}_q(2))$
\end{enumerate}
   induces a faithful representation
   $\rho_3\colon B_3 \to \mathrm{Aut}_{\mathcal{C}}(V^{\otimes 3})$
   on the fundamental representation $V$
   (\hyperref[prop:braid-action-tensor]{Proposition}).
\begin{enumerate}
  \item The image of $\rho_3$ in $\mathrm{SU}(2)$ (via the fibre
\end{enumerate}
   functor $\omega$) is the binary dihedral group
   $\mathrm{Dic}_3$ of order $12$
   (\hyperref[def:su2-generators]{Definition}).
\begin{enumerate}
  \item The composite $B_3 \xrightarrow{\;\rho_3\;}
\end{enumerate}
   \mathrm{Dic}_3 \twoheadrightarrow S_3$ recovers the three
   copies of $S_3$ inside $S_5 = \mathrm{Aut}(\Sigma_4)$
   (\hyperref[rem:chirality-trefoil]{Remark}), one for each
   colour sector.
\begin{enumerate}
  \item This structure is defined over any base ring $R$ with
\end{enumerate}
   $30 \in R^{\times}$; in particular it applies to the
   $\mathfrak{p}$-adic realisation
   (\hyperref[ex:p-adic-realisation]{Example}), where knot complements
   and skein relations are not available.

---

\emph{Item (1)} follows from \hyperref[prop:braid-action-tensor]{Proposition}
applied to the fundamental representation $V$ of
$\mathrm{SU}_q(2)$, which is $2$-dimensional.  The braiding
$c_{V,V}$ is computed by the $R$-matrix of
$\mathcal{U}_q(\mathfrak{su}(2))$
(\hyperref[rem:r-matrix-braiding]{Remark}), and faithfulness follows from the
fact that the eigenvalues of $c_{V,V}$ are $q$ and $-q^{-1}$, which
are distinct for $q \neq \pm 1$.

For \emph{item (2)}, the fibre functor $\omega$ sends the braiding to the
matrix $\check{R} = P \circ \mathcal{R}$ acting on
$\mathbb{C}^2 \otimes \mathbb{C}^2$.  The generators $\sigma_1, \sigma_2$
map to the $\mathrm{SU}(2)$ elements $\alpha, \beta$
(\hyperref[def:su2-generators]{Definition}) generating $\mathrm{Dic}_3$, as
verified in \hyperref[ex:trefoil-operator-algebra]{Example}.

\emph{Item (3)}: the projection $\mathrm{Dic}_3 \to S_3$ sends
$\alpha \mapsto (1\,2\,3)$ and $\beta \mapsto (1\,2)$.  The three
copies of $S_3$ inside $S_5$ are the stabilisers of the three pairs
of chiral coordinates
(\hyperref[def:chiral-sector-coords]{Definition}), and conjugation by
elements of $S_5 \setminus S_3$ permutes these copies.

\emph{Item (4)}: the braided monoidal structure of
$\mathrm{Rep}(\mathrm{SU}_q(2))$ is defined over $R[q, q^{-1}]$
with $30 \in R^{\times}$
(\hyperref[prop:brings-is-qou]{Proposition}(1)), and
\hyperref[prop:braid-action-tensor]{Proposition} requires no topological
input beyond the hexagon axioms.
\end{corollary}
\begin{proof}
\emph{Item (1)} follows from \hyperref[prop:braid-action-tensor]{Proposition}
applied to the fundamental representation $V$ of
$\mathrm{SU}_q(2)$, which is $2$-dimensional.  The braiding
$c_{V,V}$ is computed by the $R$-matrix of
$\mathcal{U}_q(\mathfrak{su}(2))$
(\hyperref[rem:r-matrix-braiding]{Remark}), and faithfulness follows from the
fact that the eigenvalues of $c_{V,V}$ are $q$ and $-q^{-1}$, which
are distinct for $q \neq \pm 1$.

For \emph{item (2)}, the fibre functor $\omega$ sends the braiding to the
matrix $\check{R} = P \circ \mathcal{R}$ acting on
$\mathbb{C}^2 \otimes \mathbb{C}^2$.  The generators $\sigma_1, \sigma_2$
map to the $\mathrm{SU}(2)$ elements $\alpha, \beta$
(\hyperref[def:su2-generators]{Definition}) generating $\mathrm{Dic}_3$, as
verified in \hyperref[ex:trefoil-operator-algebra]{Example}.

\emph{Item (3)}: the projection $\mathrm{Dic}_3 \to S_3$ sends
$\alpha \mapsto (1\,2\,3)$ and $\beta \mapsto (1\,2)$.  The three
copies of $S_3$ inside $S_5$ are the stabilisers of the three pairs
of chiral coordinates
(\hyperref[def:chiral-sector-coords]{Definition}), and conjugation by
elements of $S_5 \setminus S_3$ permutes these copies.

\emph{Item (4)}: the braided monoidal structure of
$\mathrm{Rep}(\mathrm{SU}_q(2))$ is defined over $R[q, q^{-1}]$
with $30 \in R^{\times}$
(\hyperref[prop:brings-is-qou]{Proposition}(1)), and
\hyperref[prop:braid-action-tensor]{Proposition} requires no topological
input beyond the hexagon axioms.
\end{proof}

\begin{remark}[Skein relations as an archimedean shadow]
  \label{rem:skein-archimedean-shadow}
  \uses{cor:descartes-braid-structure}
In the archimedean realisation, the Markov trace on the braid-group
representation $\rho_3$ yields the Jones polynomial, and the skein
relation of \hyperref[chap:experimental]{Chapter} is the trace of the
characteristic equation of $c_{V,V}$.  Over $\mathbb{Q}_\mathfrak{p}$, the
same characteristic equation holds (the eigenvalues $q$ and $-q^{-1}$
are elements of the base ring), but there is no ambient $3$-manifold in
which to interpret it as a knot invariant.  The categorical formulation
of \hyperref[cor:descartes-braid-structure]{Corollary} therefore captures
the algebraic content of braiding without presupposing an archimedean
topology.
\end{remark}

\section{The Descartes theorem}
\label{sec:descartes-theorem}

\begin{theorem}[Descartes theorem]
  \label{thm:descartes}
  \uses{def:descartes-condition, def:brings-surface, prop:brings-is-qou, thm:qu-to-qou, prop:genus-maximal-descent}
Let $\mathcal{C}_{-6}$ denote the category of quantum observable universes
over $R$ satisfying the Descartes condition $\chi(A) = -6$
(\hyperref[def:descartes-condition]{Definition}), with morphisms given by QOU maps
that preserve the Euler characteristic.  Then Bring's surface $\Sigma_4$
(\hyperref[def:brings-surface]{Definition},
\hyperref[prop:brings-is-qou]{Proposition}) is a terminal object in
$\mathcal{C}_{-6}$: for every Descartes universe
$(\mathcal{C}, A, \omega)$ in $\mathcal{C}_{-6}$, there exists a unique
morphism to the Bring's QOU.

---

\emph{Existence.}
By \hyperref[thm:qu-to-qou]{Theorem} and
\hyperref[prop:genus-maximal-descent]{Proposition}, the base curve is a smooth
genus-$4$ curve $\Sigma$ with $\chi(\Sigma) = -6$.
The Hurwitz bound gives $\lvert\mathrm{Aut}(\Sigma)\rvert \leq 252$;
among genus-$4$ curves, Bring's surface achieves the maximum
$\lvert S_5\rvert = 120$ (Edge, 1978), and the $S_5$-equivariant structure
uniquely determines $\Sigma_4 \subset \mathbb{P}^4_R$.
The descent data (branch locus with $6$ points, $S_5$-symmetric
ramification) defines a map $\Sigma \to \Sigma_4$, extending to a QOU
morphism by functoriality of the Haar state on
$\mathcal{O}_q(\mathrm{SU}(2))$.

\emph{Uniqueness.}
By Tannaka--Krein rigidity, any two fibre-preserving monoidal functors
between Tannakian categories are naturally isomorphic.
\end{theorem}
\begin{proof}
\emph{Existence.}
By \hyperref[thm:qu-to-qou]{Theorem} and
\hyperref[prop:genus-maximal-descent]{Proposition}, the base curve is a smooth
genus-$4$ curve $\Sigma$ with $\chi(\Sigma) = -6$.
The Hurwitz bound gives $\lvert\mathrm{Aut}(\Sigma)\rvert \leq 252$;
among genus-$4$ curves, Bring's surface achieves the maximum
$\lvert S_5\rvert = 120$ (Edge, 1978), and the $S_5$-equivariant structure
uniquely determines $\Sigma_4 \subset \mathbb{P}^4_R$.
The descent data (branch locus with $6$ points, $S_5$-symmetric
ramification) defines a map $\Sigma \to \Sigma_4$, extending to a QOU
morphism by functoriality of the Haar state on
$\mathcal{O}_q(\mathrm{SU}(2))$.

\emph{Uniqueness.}
By Tannaka--Krein rigidity, any two fibre-preserving monoidal functors
between Tannakian categories are naturally isomorphic.
\end{proof}

\begin{remark}[Universality for a given Euler characteristic]
  \label{rem:universality-general}
  \uses{thm:descartes}
The Descartes theorem (\hyperref[thm:descartes]{Theorem}) is expected to be a
special case of a more general result: for any admissible Euler
characteristic $\chi$, the category $\mathcal{C}_{\chi}$ of $\mathbb{C}$
quantum observable universes with $\chi(E) = \chi$ should admit a terminal
object determined by the unique maximally symmetric Riemann surface of the
corresponding genus.  The Descartes case $\chi = -6$ (genus $4$) is
distinguished by Bring's curve, but analogous universal objects are expected
at other values of $\chi$.
\end{remark}

\begin{example}[Bring's surface and Yekutieli obstructions]
  \label{ex:descartes-yekutieli}
  \uses{thm:qu-to-qou, rem:mass-geometric-drag, rem:chirality-trefoil}
Bring's surface satisfies \hyperref[thm:qu-to-qou]{Theorem} with
$\chi(A) = -6$.  The non-vanishing Yekutieli classes $\kappa_n \neq 0$
(\hyperref[thm:qu-to-qou]{Theorem}(3)) obstruct smooth lifting at every level,
creating twisted sectors that control inertial mass
(\hyperref[rem:mass-geometric-drag]{Remark}) and the three lepton generations
--- corresponding to the three pairs of branch points under
$S_3 \subset S_5$ (\hyperref[rem:chirality-trefoil]{Remark}).
\end{example}

\begin{remark}[Observations on the Bring surface]
  \label{rem:observation-brings}
  \uses{def:observation, def:substrate-parameter, rem:two-observer-parallelogram}
The general notion of an observation
(\hyperref[def:observation]{Definition}) specialises on the Bring surface as
follows: the observation morphism $\mathcal{O} \colon A \to A$ acts
fibrewise on the state bundle $E \to M$ and perturbs the substrate
parameter $q(p)$ (\hyperref[def:substrate-parameter]{Definition}), shifting
the non-commutative geometry of the quantum Bring surface
(\hyperref[sec:quantum-brings]{Section}).  The two-observer parallelogram
(\hyperref[rem:two-observer-parallelogram]{Remark}) applies directly: two
independent observations of a fermion on the Bring surface produce a
discrepancy controlled by the local Planck constant $\hbar_q(p)$.
\end{remark}
